\chapter{Introduction}
\label{ch:introduction}

\section{Motivation}

Deploying a neural network outside a data centre means making it smaller. The
standard tool is quantization: representing weights and activations with fewer
bits, trading numerical precision for memory and energy. The operation is
routinely validated by checking that accuracy does not fall, and for many
applications that is a sufficient test.

It is not sufficient when the model is used to explain something. Rabbi et
al.~\cite{rabbi2026} evaluated five convolutional architectures at INT8 and
INT4, comparing the attribution maps of the quantized models against their
full-precision counterparts. They found that explanation similarity degrades
measurably between INT8 and INT4 even where accuracy differences are negligible,
and that when quantization does cause a misclassification, the quantized model
is generally attending to \emph{different} regions of the input rather than
weighing the same evidence differently. Accuracy and interpretability, in other
words, do not degrade together.

That finding is confined to conventional artificial neural networks. Spiking
neural networks, the architecture most closely associated with low-power
deployment, differ from them in ways that make the question both harder and more
interesting.

\section{Research question}

The question this thesis addresses is:

\begin{quote}
\emph{Does a quantized spiking neural network under edge-realistic constraints
base its decisions on the same evidence as its full-precision counterpart, and
at what energy cost?}
\end{quote}

Three properties of spiking networks have no counterpart in the convolutional
setting, and each one bears on the question.

\paragraph{A second quantization target.} A spiking neuron carries a membrane
potential: a state variable read and written at every time step, and the
dominant memory cost at deployment. It is quantized independently of the
weights~\cite{yin2024mint}, whereas a conventional network offers only weights
and activations. The two targets can therefore be separated, which converts a
single observation into a causal ablation.

\paragraph{A temporal axis in the explanation.} A spiking explanation is a
\emph{sequence} of maps, one per time step. Quantization may preserve where a
network looks while displacing when it decides, a change to which spatial
metrics are blind by construction.

\paragraph{A native explainer.} The Spike Activation
Map~\cite{kim2021sam} reads spike activity directly. It requires neither
gradients, side-stepping the approximation error a surrogate gradient injects
into any transplanted Grad-CAM-style method~\cite{selvaraju2017gradcam}, nor
ground-truth labels.

The working hypothesis follows from the first two properties: \textbf{quantizing
the membrane potential degrades the temporal structure of the explanation before
quantizing the weights degrades its spatial structure}, since the membrane
potential is the variable that integrates evidence over time. The hypothesis is
falsifiable, and testing it requires quantizing the two targets separately.

\section{What this thesis establishes}

Answering the research question requires an instrument, and building the
instrument turned out to be the substantive part of the work. The metrics used
to compare explanations return numbers between zero and one whatever is fed to
them; before a quantization result can be read, one must know what they report
when nothing has changed.

This thesis therefore contributes the following.

\begin{enumerate}
  \item \textbf{A hardware-realizable full-precision baseline}, with every
        design choice settled on evidence and recorded in a frozen manifest: the
        input resolution chosen on network health rather than accuracy, the
        membrane decay restricted to shift-implementable values inside the
        hyperparameter search rather than rounded afterwards, and a three-seed
        reference at $0.9833 \pm 0.0014$ validation accuracy.

  \item \textbf{A corrected and re-derived implementation of the Spike
        Activation Map}, with its decay parameter fixed on the temporal axis it
        actually governs, and its faithfulness tested quantitatively against a
        random baseline rather than judged by visual plausibility.

  \item \textbf{A calibration of the explanation-divergence metrics} at three
        levels: numerical noise, a thirtyfold change in the explainer, and
        initialization variance. This establishes what those metrics can and
        cannot detect, and it is the contribution with the widest scope, since
        neither of the two closest works~\cite{rabbi2026, smith2026} establishes
        any of the three.

  \item \textbf{A complete specification of the quantization study} the
        calibration was built to support, together with the limits the
        calibration imposes on what that study will be able to claim.
\end{enumerate}

The central result is negative in form and useful in consequence. Two
full-precision models differing only in random seed, and indistinguishable in
accuracy, produce Spike Activation Maps correlating at $0.153$ on the worst
class. Spatial explanation-similarity metrics cannot separate a quantization
effect from initialization variance in this architecture. The temporal centre of
mass, being an aggregate over the whole map rather than a pixel-wise comparison,
survives the same test and becomes the primary instrument.

\section{Application domain}

The classification task is the identification of transient noise --- known as
\emph{glitches} --- in the data of the LIGO gravitational-wave detectors, using
the Gravity Spy release~\cite{zevin2017, glanzer2023, glanzer2021}.

The domain was chosen for three reasons. Explanation fidelity has an
\emph{operational} consequence there: glitch classification guides detector
characterization, where a commissioner acts physically on the instrument, so a
compressed model reaching the same label from different evidence would preserve
the label while corrupting the diagnostic cue. The classes have physically
interpretable morphology, so a domain expert can judge whether an attention map
falls where the glitch actually is, which is not true of a generic image
benchmark. And the data is public, large and labelled.

\section{Scope}

This work characterizes explanation fidelity under \emph{simulated}
quantization, with analytical energy and memory accounting. Hardware deployment
is explicitly out of scope: no model is exported to a neuromorphic device and no
latency is measured on one.

Every design decision in this thesis is taken on the validation partition. The
test partition is opened once, for the final estimate reported in
Section~\ref{sec:test}.

\section{Structure}

Chapter~\ref{ch:background} reviews spiking neural networks, their training by
surrogate gradients, quantization for spiking architectures, explainability
methods applicable to them, and the gravitational-wave classification task.
Chapter~\ref{ch:methodology} describes the data pipeline, the model, the
selection of the operating configuration and the implementation of the Spike
Activation Map. Chapter~\ref{ch:results} reports the full-precision reference,
the physical validation of the network dynamics, the explainability analyses and
the calibration of the divergence metrics. Chapter~\ref{ch:conclusions}
discusses the limits of these results and specifies the quantization study that
follows.
